\section{算子理论}\label{chap:operator-theory}

本章简要讨论从 Banach 空间 $X$ 到 Banach 空间 $Y$ 的线性算子 $A$ 的理论. 主要关注方程
\begin{equation}\label{eq:abstract-linear-operator-equation}
Ax=y,
\end{equation}
其中给定 $y\in Y$, 所讨论的主要问题是解 $x\in X$ 的存在性, 多重性和可计算性. 当然, 大多数读者在线性代数的学习中已经讨论过这些问题. 在那里, 空间 $X$ 和 $Y$ 分别为有限维向量空间 $\mathbb{R}^n$ 和 $\mathbb{R}^m$, $A$ 由一个 $m\times n$ 矩阵表示. 在本书的 Definition~\ref{def:bounded-linear-mapping} 中, 定义了从一个 (可能无穷维的) Banach 空间到另一个 Banach 空间的有界线性算子, 已经考察过一种更一般类型的算子. 然而, 正如下面将看到的, PDE (以及 ODE) 中许多重要算子无界. 强烈建议读者将本节结果与以前的线性代数教材中的结果比较, 同时记住理论的两项主要推广: 推广到无穷维空间, 以及推广到无界算子.

\textbf{Note:} 虽然将算子定义为 Banach 空间之间的映射, 但本书所讨论的算子理论应用, 大多涉及可分 Hilbert 空间之间的映射. 因而在下面的许多定理中, 只对 Hilbert 空间或可分 Hilbert 空间的情形给出陈述或证明. 这种做法大大减少了需要发展的工具, 但也限制了仅用本书材料所能处理的应用范围. 这是在 ``街头'' 学习泛函分析所付出的代价之一.

下文中, 除非另有说明, 将用记号 $X$ 和 $Y$ 表示 Banach 空间, 用 $H$ 表示 Hilbert 空间.

\subsection{基本定义与例子}\label{sec:operator-basic-definitions-examples}

\subsubsection{算子}\label{subsec:operators}

为容纳无界算子, 本节先给出以下扩充的定义.

\begin{Definition}
\label{def:linear-operator-with-specified-domain}
设 $X$ 和 $Y$ 为 Banach 空间. 从 $X$ 到 $Y$ 的线性算子是一个对 $(\mathcal{D}(A),A)$, 由子空间 $\mathcal{D}(A)\subset X$ (称为算子的\textbf{定义域}) 和线性变换 $A:\mathcal{D}(A)\to Y$ 组成.
\end{Definition}

许多数学学生都不得不忍受这样的微积分教师: 坚持认为, 定义域为 $[0,1]$ 的函数 $f(x)=x$ 与定义在整个实直线上的同一函数有根本区别. 学生很快意识到, 在大多数情况下可以忽略这种区别. 在本章中将看到, 将定义域纳入算子的定义并非仅仅吹毛求疵. 对无界算子, 指明定义域确实可能造成实质区别. 然而, 如此强调定义域在算子定义中的重要性之后, 又常会使用忽略这一点的不严谨语言. 即经常只说 ``算子 $A$'', 而不指明定义域. 这种用法在有界算子的研究中是标准且无歧义的 (将在下面的 Theorem~\ref{thm:bounded-operator-norm-preserving-extension} 中看到, 其定义域可以延拓到整个 $X$), 在没有混淆可能时, 即使对无界算子也沿用这种较简短的称谓.

将同时使用记号 $Ax$ 和 $A(x)$ 表示算子对其定义域中元素的作用.

\begin{Definition}
\label{def:operator-range-null-space}
$(\mathcal{D}(A),A)$ 的值域是由下式定义的子空间 $\mathcal{R}(A)\subset Y$:
\begin{equation}\label{eq:linear-operator-range-definition}
\mathcal{R}(A):=\{u\in Y\mid u=A(x),\quad\text{for some }x\in\mathcal{D}(A)\}.
\end{equation}
$(\mathcal{D}(A),A)$ 的零空间是由下式定义的子空间 $\mathcal{N}(A)\subset X$:
\begin{equation}\label{eq:linear-operator-null-space-definition}
\mathcal{N}(A):=\{x\in X\mid A(x)=0\}.
\end{equation}

\end{Definition}

如此定义值域之后, 可对算子 $(\mathcal{D}(A),A)$ 使用以下记号:\[X\supset\mathcal{D}(A)\ni x\mapsto A(x)\in\mathcal{R}(A)\subset Y.\]

有时将集合 $X$ 和 $Y$ 分别称为\emph{余值域}和\emph{陪域}, 以区别于它们各自的子空间, 即定义域和值域. 虽然认同这种区分的重要性, 但不采用这些术语.
